\section{Language Hypothesis Testbench}

\subsection{Purpose}

This section tests language and script hypotheses against the current structural, abstract phonetic, and semantic models. It does not claim decipherment. Its job is to separate three levels of claim:

\begin{itemize}
  \item a script-function claim, such as restricted administrative or onomastic writing;
  \item a language-family compatibility claim, such as Dravidian-compatible suffixing structure;
  \item a true language identification claim, which would require validated phonetic readings.
\end{itemize}

The project is now strong enough to score the first two levels. It is not yet strong enough to make the third.

\subsection{Method}

The script \texttt{scripts/language-hypothesis-testbench.py} converts each hypothesis into explicit diagnostics and scores those diagnostics against existing outputs:

\begin{itemize}
  \item structural formulae and semantic frames;
  \item abstract phonetic variables and minimal-pair evidence;
  \item semantic context associations with site, region, artifact type, material, and iconography;
  \item strong onomastic anchors in the \texttt{002 X 740} prime-entity frame.
\end{itemize}

The model penalizes all language-family hypotheses because no sign has yet received a validated sound value. This is intentional. Compatibility is not identification.

\subsection{Current Result}

\begin{table}[htbp]
\centering
\begin{tabular}{lp{0.46\textwidth}l}
\toprule
\textbf{Rank} & \textbf{Hypothesis} & \textbf{Score} \\
\midrule
1 & Restricted administrative/onomastic script & 0.731 \\
2 & Multilingual administrative ecology & 0.689 \\
3 & Primarily non-linguistic emblem system & 0.621 \\
4 & Dravidian-compatible linguistic layer & 0.605 \\
5 & Logo-syllabic or logo-phonetic writing & 0.585 \\
6 & Indo-Iranian/Indo-Aryan linguistic layer & 0.541 \\
7 & Munda/Para-Munda linguistic layer & 0.527 \\
\bottomrule
\end{tabular}
\caption{Hypothesis ranking. Scores are compatibility scores, not proof.}
\label{tab:language-hypothesis-ranking}
\end{table}

The best overall model is currently a restricted administrative/onomastic script. The best language-family-compatible model is Dravidian-compatible, mostly because the corpus now shows strong suffix/final behavior and productive name/title-like slots. However, this is not enough to identify the language. Indo-Iranian/Indo-Aryan and Munda/Para-Munda remain plausible but less strongly supported by the current structural evidence.

\subsection{Diagnostic Features}

\begin{table}[htbp]
\centering
\begin{tabular}{lr}
\toprule
\textbf{Feature} & \textbf{Score} \\
\midrule
Formula restriction & 0.898 \\
Context association & 0.890 \\
Usable iconography & 0.794 \\
Suffix/final structure & 0.734 \\
Regional differentiation & 0.733 \\
Prime entity structure & 0.654 \\
Phonetic readiness & 0.583 \\
Classifier/title structure & 0.517 \\
\bottomrule
\end{tabular}
\caption{Diagnostic evidence features.}
\label{tab:language-diagnostic-features}
\end{table}

The suffix/final score is the strongest reason Dravidian remains interesting. The classifier/title and context-association scores are the strongest reasons a language-neutral or multilingual administrative model remains stronger than any single-language claim.

\subsection{What We Can Claim Now}

\begin{itemize}
  \item The inscriptions are highly compatible with restricted formulaic writing.
  \item The \texttt{002 X 740} frame is the best current partial-decipherment target.
  \item \texttt{X} is likely an entity slot: name, title, office, lineage, place, institution, commodity, or authority formula.
  \item Dravidian-compatible suffixing is currently the strongest language-family-compatible pattern.
  \item No language-family identification is justified until at least one phonetic or lexical value predicts all occurrences of a strong anchor.
\end{itemize}

\subsection{Lexical Reading Gate}

The file \texttt{outputs/lexical\_reading\_gate.csv} defines what kinds of lexical searches are allowed for each unit. For the strongest prime-entity anchors, lexical searches are limited to proper names, offices, titles, lineages, places, institutions, and emblem-linked identities. The model explicitly blocks shortcuts such as equating a sign with an animal name just because the seal has that animal.

The immediate reading gate for the top anchors is:

\begin{itemize}
  \item \texttt{176}: search only after all Mohenjo-daro and Lothal occurrences are image-checked.
  \item \texttt{235-222}: test bull-linked name/title readings, but only if the Chanhu-daro cluster holds visually.
  \item \texttt{032-220}: test cross-site emblem/name readings across bull-seal contexts.
  \item \texttt{240-482}: test Harappa administrative or institutional readings before personal-name readings.
  \item \texttt{590-390}: test steatite seal name/title/place readings after image validation.
\end{itemize}

\subsection{Interpretation}

The project has not cracked the language yet, but it has narrowed the problem. The current best model is:

\begin{center}
\texttt{RESTRICTED ADMINISTRATIVE/ONOMASTIC SCRIPT}
\end{center}

with a probable formula:

\begin{center}
\texttt{CLASSIFIER/TITLE + ENTITY + TERMINAL/SUFFIX}
\end{center}

and a prime target:

\begin{center}
\texttt{002 + ENTITY + 740}
\end{center}

The strongest next decipherment route is therefore not a global alphabet or syllabary. It is controlled proper-name and title testing, using image-grounded semantic anchors and minimal alternations.

\subsection{Outputs}

\begin{itemize}
  \item \texttt{outputs/language\_evidence\_features.csv}
  \item \texttt{outputs/language\_hypothesis\_scores.csv}
  \item \texttt{outputs/lexical\_reading\_gate.csv}
  \item \texttt{outputs/decipherment\_claims.csv}
  \item \texttt{outputs/language\_hypothesis\_summary.csv}
  \item \texttt{outputs/language\_hypothesis\_testbench.tex}
\end{itemize}
